\section{One Attribute, and Three Subgraphs in the Error}
\label{sec:ch15-one-attribute-and-three-errors}

\index{resolvability!the mistake this chapter opens on|(}%
Chapter~\ref{ch:null-blast-radius} ended on a graph whose failures are decided
by the schema, with no operational lever anywhere near them. This chapter is
about the step that is supposed to check that schema, and it starts by getting
one wrong.

\index{resolvable@\code{resolvable: false}!the argument that reads like modesty}%
Chapter~\ref{ch:composition} said in one sentence that \code{resolvable: false}
is the right thing to say about a type you only point at and the wrong thing to
say about a type you contribute fields to.
Chapter~\ref{ch:second-third-subgraph} then used it correctly, on the
\code{Speaker} stub: the Sessions service names the type so that
\code{Session.speaker} has one, declares which field identifies a speaker, and
writes \code{[Key("id", false)]} to say it cannot turn such a key back into
anything. One sentence and one correct use is not much to hold an argument that
reads like modesty. Read in English it says \emph{this service does not own
these rows}, which is true of the Sessions service and true of the next service
a reader will apply it to, where it is also wrong.

\index{Ratings subgraph!does not own Session rows either}%
The Ratings service does not own \code{Session} rows either. It holds scores,
and everything interesting it knows is about a type another team maintains. So
put the same argument on its key, in \code{SessionType.cs}, and take the
\code{[ReferenceResolver]} attribute off with it, because a key nothing may
enter through has nothing to route to a resolver.

\index{Ratings subgraph!what the exported schema loses}%
The exported schema changes in two places rather than one. The key gains its
argument, and \code{\_entities} and the \code{\_Entity} union disappear from the
document entirely, because a service nothing may enter has no entity route to
publish. Chapter~\ref{ch:second-third-subgraph} recorded that from the other
side, on the stub, where it was the point rather than a surprise.

Compose the graph.

\subsection{Three Errors, and Where They Point}

\index{composition!one error per unresolvable field}%
It fails three times, once for each field Ratings contributes to a
\code{Session} it can no longer be asked about. The three blocks are identical
apart from the field name. Here is the first, with the box \code{wgc} draws
around it stripped:

\begin{minted}[fontsize=\footnotesize,breakanywhere]{text}
The field "averageScore" is unresolvable at the following path:
 query {
  searchSessions {
   edges {
    node {
     session {
      averageScore <--
     }
    }
   }
  }
 }
This is because:
 - The root type field "Query.searchSessions" is defined in the following subgraph: "search".
 - The field "Session.averageScore" is defined in the following subgraph: "ratings".
 - The entity ancestor "Session" in subgraph "sessions" has no accessible target entities (resolvable @key directives)
in the subgraphs where "Session.averageScore" is defined.
 - The type "Session" is not a descendant of any other entity ancestors that can provide a shared route to access
"averageScore".
\end{minted}

\index{composition!three subgraphs in one error}%
Count the subgraphs in that message. The path starts in \code{search}. The
entity ancestor is in \code{sessions}. The field is in \code{ratings}. Three
services, and I changed one line in the third of them.

\index{composition!where the fix actually is}%
The line that names the file to open is the second of the four, and it is the
least dramatic thing on the page. Everything above it describes a query nobody
sent, through a service that has nothing wrong with it, and the arrow at the
end of that query points at a field sitting exactly where it belongs.

\index{composition!reading it from the bottom}%
The path is the composer showing its working, and the working is not the
diagnosis. The second reason line is the one that names a type, a field and the
subgraph that declares them, and that subgraph is the file to open.
Chapter~\ref{ch:composition} pointed at the third line instead, and both are
right about different things: the third says what is wrong in the composer's
vocabulary, and the second says where.
Section~\ref{sec:ch15-reading-one-properly} is the rest of the procedure.

\index{Search subgraph!a bystander in somebody else's error}%
Chapter~\ref{ch:cross-seam-paging} built the Search service, and it is in this
message because its discovery documents carry a \code{session} field that
reaches the type. It is a bystander. It has no opinion about ratings, it does
not know the Ratings service exists, and taking it out of the graph changes
what this message says without changing one character of the mistake. The next
section is about how much it changes.
\index{resolvability!the mistake this chapter opens on|)}%
